%% 05_devices.tex -- Chapter 5: Devices for Microwatt Design
%% Status: OUTLINE (build stage 3). Replace the outline with the full chapter when it is written.
\chapter{Devices for Microwatt Design}\label{ch:devices}
\chapterlead{A biopotential amplifier lives on a few hundred nanoamperes per transistor, so its devices sit in weak or moderate inversion, where the square law no longer applies.}

\begin{specnote}[Chapter status: outline]
This chapter is an outline. It will be written in build stage 3 (Part II with Labs 3--5). The topics, the design question and the labs below are fixed; the derivations, worked examples and figures are still to come.
\end{specnote}

\begin{objectives}
\begin{itemize}
  \item place a MOSFET in weak, moderate or strong inversion from its inversion coefficient and read its $\gm/I_\mathrm{D}$;
  \item explain why $\gm/I_\mathrm{D}$ peaks at about $1/(n U_T)$ in weak inversion and what that means for noise per microampere;
  \item relate flicker noise to gate area and pick input-pair dimensions for a sub-\qty{100}{\hertz} band;
  \item compare \sgx{} MOSFETs and SiGe HBTs as low-noise input devices (thermal noise, flicker noise, base current);
  \item estimate offset and gain mismatch from Pelgrom coefficients.
\end{itemize}
\end{objectives}

\begin{insight}[Design question]
Which device, at what inversion level and size, gives the lowest input-referred noise for a microwatt bias current?
\end{insight}

\section{Weak, moderate and strong inversion}\label{sec:ch5-1}
The inversion coefficient $IC = I_\mathrm{D}/I_\mathrm{spec}$ places a transistor on one continuous curve from exponential (weak inversion) to square-law (strong inversion) behaviour. Biomedical front-ends bias their input pairs at $IC \lesssim 1$.

\section{The \texorpdfstring{$\gm/I_\mathrm{D}$}{gm/ID} method at microwatt bias}\label{sec:ch5-2}
Sizing from $\gm/I_\mathrm{D}$ lookup tables instead of hand equations; the weak-inversion limit $\gm/I_\mathrm{D} \approx 1/(nU_T) \approx \qty{25}{\per\volt}$; why it saturates and what it costs in area and speed. This chapter cross-links \emph{The Art of Analog IC Design}~\cite{artofaicd} rather than repeating it.

\section{Flicker noise and device area}\label{sec:ch5-3}
$1/f$ noise scales as $1/(WL)$, so the low corner of an EEG band forces large input devices; PMOS versus NMOS in \sgx{}; the corner frequency and its dependence on bias.

\section{MOSFET versus SiGe HBT input pairs}\label{sec:ch5-4}
An HBT reaches $\gm = I_\mathrm{C}/U_T$ at any current and has very low $1/f$ noise, but its base current flows through the electrode impedance and adds shot noise. The trade-off decides the input stage of the \OB{} LNA.

\section{Matching}\label{sec:ch5-5}
Pelgrom's law $\sigma(\Delta V_\mathrm{T}) = A_{V_T}/\sqrt{WL}$, current-mirror and capacitor matching, and why large, slow biomedical devices match well but still need chopping for microvolt offsets.

\begin{tryit}[Labs for this chapter]
Lab~3 (\cref{sec:lab3}) characterises $\gm/I_\mathrm{D}$, flicker noise versus $W\!\cdot\!L$ and HBT versus MOS input noise in \sgx{}; its lookup tables are reused in Labs~4 and~5.
\end{tryit}

\paragraph{Read with this chapter} Sarpeshkar~\cite{sarpeshkar} (subthreshold design, noise--power limits), Razavi~\cite{razavi} (device noise), \emph{The Art of Analog IC Design}~\cite{artofaicd} (Part II: the transistor and $\gm/I_\mathrm{D}$).
